\chapter[Introdução]{Introdução}

Jogos digitais competitivos constituem hoje um dos segmentos mais expressivos da indústria do entretenimento, sustentados por cenas profissionais consolidadas e por comunidades que avaliam o produto sob critérios técnicos rigorosos \cite{yannakakis2018ai}. Entre os subgêneros desse universo, os jogos de luta --- também referidos pela sigla FGC, do inglês \textit{Fighting Game Community} --- ocupam posição particular: duelos sucessivos entre dois personagens, em tempo real, em que cada participante escolhe um lutador dentre um elenco com perfis funcionais distintos. O desfecho de cada confronto depende menos da reação imediata do que do casamento estrutural entre as capacidades dos dois personagens escolhidos, o que torna o desenho desse elenco o problema central de \textit{design} do subgênero.

O conceito que organiza esse problema é o de \textbf{equilíbrio competitivo}: a propriedade de que nenhum personagem do elenco apresente vantagem suficientemente acentuada sobre os demais a ponto de tornar a escolha trivial ou de inviabilizar competitivamente os personagens preteridos \cite{adams2014fundamentals,elias2012characteristics}. Em jogos de luta, o equilíbrio raramente se manifesta como simetria estatística pura entre personagens; manifesta-se, com mais frequência, como uma estrutura de vantagens recíprocas e não-transitivas, em que cada personagem vence alguns confrontos e perde outros, mas nenhum domina o conjunto. \citeonline{gingold2006rps} formalizou essa propriedade ao demonstrar, por construção, que jogos de luta competitivos comportam-se como variantes em tempo real do clássico Pedra-Papel-Tesoura, herdando dele a estrutura cíclica de vantagens como fundamento estrutural do gênero.

Tradicionalmente, o trabalho de produzir um elenco equilibrado é conduzido de forma manual por equipes de \textit{design}, que ajustam iterativamente os atributos numéricos de cada personagem, observam o comportamento emergente em partidas controladas e refinam os valores até que nenhum personagem se imponha sobre os demais. O processo é caro, lento e propenso a erros: requer múltiplos ciclos de teste, depende de jogadores habilitados a estressar o sistema sob condições competitivas, e frequentemente exige correções pós-lançamento após a comunidade competitiva identificar configurações dominantes que escaparam ao processo interno. A literatura em computação evolutiva tem investigado, há ao menos duas décadas, a aplicação de algoritmos baseados em busca como ferramenta para automatizar parte desse esforço, em domínios que vão de jogos de tabuleiro \cite{hom2007automatic,browne2010evolutionary} a \textit{role-playing games} massivos \cite{chen2014mmorpg}, passando por geração procedural de conteúdo em sentido mais amplo \cite{togelius2011searchbased}.

No entanto, há uma tensão própria ao problema do balanceamento que a literatura aborda apenas parcialmente. Equilibrar matematicamente um elenco é trivial no limite: atribuir a todos os personagens o mesmo conjunto de atributos garante que nenhum vença sistematicamente os demais. Esse equilíbrio degenerado, porém, destrói o elemento que distingue o subgênero --- a existência de personagens com identidades funcionais reconhecíveis, cada um associado a um modo característico de jogar. A pressão por equilíbrio competitivo, quando aplicada sem contrapeso, tende a homogeneizar o elenco, convergindo todos os personagens para perfis genéricos e neutros que perdem a singularidade que justifica o próprio elenco. O problema que motiva este trabalho é, portanto, a tensão entre duas dimensões legítimas e simultaneamente desejáveis: o equilíbrio competitivo do elenco e a preservação das identidades funcionais que distinguem seus personagens.

\section{Hipótese}

Este trabalho parte da seguinte hipótese: \textit{é possível, mediante a aplicação de Algoritmos Genéticos multi-objetivo, obter equilíbrio competitivo entre arquétipos distintos de personagens em jogos de luta sem destruir suas identidades funcionais; e o trade-off entre essas duas dimensões pode ser tornado explícito por meio da fronteira de Pareto obtida pelo NSGA-II} \cite{deb2002nsga}.

A hipótese pressupõe que a preservação de identidade arquetípica seja \textbf{mensurável} de forma quantitativa e independente do objetivo de equilíbrio. Esse pressuposto distingue a presente investigação de abordagens que tratam o balanceamento exclusivamente como problema escalar de equalização de taxas de vitória, e situa o trabalho na intersecção entre balanceamento competitivo e preservação de diversidade funcional, intersecção pouco explorada na literatura de busca aplicada a jogos.

\section{Justificativa}

A relevância acadêmica deste trabalho situa-se em três frentes complementares.

A primeira é a delimitação do problema. A literatura de balanceamento de jogos via algoritmos evolutivos concentra-se predominantemente em domínios distintos do recorte adotado aqui. \citeonline{hom2007automatic} aplicaram Algoritmos Genéticos ao balanceamento de jogos de tabuleiro, e \citeonline{browne2010evolutionary} estenderam a abordagem para a geração evolutiva de regras de jogos. Mais recentemente, \citeonline{chen2014mmorpg} aplicaram programação coevolutiva ao balanceamento das funções de crescimento de atributos por raça em \textit{Massively Multiplayer Online Role-Playing Games} de combate por turnos, abordagem que constitui o trabalho mais próximo do problema aqui tratado em termos de objeto (balanceamento entre classes de personagens em PvP). Nenhum desses trabalhos, todavia, opera sobre jogos de luta em tempo real --- domínio cujas mecânicas de combate, dependentes de posicionamento contínuo e de janelas de ataque sub-segundo, impõem requisitos distintos sobre o modelo de simulação --- nem aborda a preservação de identidade arquetípica como objetivo formalmente medido e otimizado.

A segunda frente é metodológica. Tratar simultaneamente o equilíbrio competitivo e a preservação de identidade requer formular uma função de aptidão capaz de representar ambas as dimensões sem que uma colapse a outra. Este trabalho propõe uma formulação multi-componente em que a dominância competitiva entre personagens e o desvio em relação a perfis arquetípicos canônicos são medidos separadamente, e investiga duas estratégias de otimização sobre essa formulação: um Algoritmo Genético escalar, que pondera as componentes em uma única função de aptidão, e o NSGA-II \cite{deb2002nsga}, que preserva o trade-off entre as componentes ao longo de toda a evolução, expondo-o de forma explícita por meio da fronteira de Pareto. A comparação entre as duas estratégias sobre o mesmo motor de simulação caracteriza, em si, uma contribuição metodológica.

A terceira frente é a postura epistemológica do trabalho. A estrutura cíclica de vantagens entre arquétipos --- na qual cada arquétipo vence alguns confrontos e perde outros --- não é codificada no processo evolutivo como restrição. Os Algoritmos Genéticos atuam livremente sobre o espaço de atributos, e a preservação da estrutura cíclica é avaliada \textit{a posteriori} como métrica de resultado, jamais como pressão evolutiva. Essa escolha é deliberada: codificar a estrutura cíclica como restrição tornaria a investigação circular, pois o sistema preservaria a estrutura porque foi pago para preservá-la, e não porque ela emerge das mecânicas e da pressão por equilíbrio. Trata-se, portanto, de uma investigação sobre o que o algoritmo é capaz de descobrir quando se lhe oferece o mínimo de viés direcional.

\section{Objetivos}

\subsection{Objetivo geral}

Investigar a aplicabilidade de Algoritmos Genéticos multi-objetivo ao problema do balanceamento competitivo de personagens em jogos de luta, considerando explicitamente a preservação de identidade arquetípica como objetivo mensurável e ortogonal ao equilíbrio.

\subsection{Objetivos específicos}

\begin{enumerate}
    \item Modelar cinco arquétipos canônicos da convenção de jogos de luta --- \textit{Rushdown}, \textit{Zoner}, \textit{Grappler}, \textit{Combo Master} e \textit{Turtle} --- como vetores numéricos de atributos físicos e pesos comportamentais.
    \item Implementar um simulador de combate um-contra-um determinístico no cômputo de dano, em que a estocasticidade restringe-se à política comportamental do agente, de modo a maximizar a sensibilidade do processo evolutivo a perturbações genéticas.
    \item Formular uma função de aptidão multi-componente que represente separadamente o equilíbrio competitivo do elenco --- medido como dominância entre personagens em confrontos diretos --- e a preservação de identidade arquetípica --- medida como desvio em relação aos perfis canônicos --- sem codificar o ciclo de vantagens entre arquétipos como restrição.
    \item Aplicar duas estratégias de otimização sobre a mesma formulação --- um Algoritmo Genético escalar e o NSGA-II \cite{deb2002nsga} --- e analisar comparativamente os elencos evoluídos por cada estratégia.
    \item Validar a preservação de identidade arquetípica nos elencos evoluídos por meio de asserções estruturais sobre rankings ordinais de atributos, complementarmente à métrica contínua de desvio utilizada na função de aptidão.
\end{enumerate}

\section{Estrutura do trabalho}

O restante deste trabalho está organizado em seis capítulos. O \autoref{chap: fundamentacao} apresenta a fundamentação teórica, revisando os conceitos de Algoritmos Genéticos, otimização multi-objetivo via NSGA-II, balanceamento em jogos competitivos e a noção de arquétipos em jogos de luta. O \autoref{chap: metodologia} descreve a metodologia adotada, incluindo a modelagem dos cinco arquétipos, a especificação do simulador de combate, a formulação da função de aptidão multi-componente e a configuração dos dois algoritmos evolutivos. O \autoref{chap: implementacao} detalha a implementação computacional do sistema, com ênfase nas decisões de engenharia que sustentam a sensibilidade do processo evolutivo. O \autoref{chap: resultados} reporta os resultados experimentais obtidos pelas duas estratégias evolutivas, complementados por análises de sensibilidade dos atributos e por validação estrutural das identidades arquetípicas. O \autoref{chap: discussao} discute os achados à luz da hipótese e das limitações do modelo. O \autoref{chap: conclusao} apresenta as conclusões do trabalho e indica direções para investigações futuras.